\section{Related work}
\label{sec:related}

\subsection{Classical Camera Calibration with Geometric Primitives}

Camera calibration remains a foundational problem in computer vision. Traditionally, it has been addressed by observing a set of well-known 3D features such as checkerboards, points, lines, or circles projected into the image plane. These features provide algebraic constraints between the 3D world and the 2D image allowing the estimation of intrinsic and extrinsic camera parameters. While effective, these approaches require specialized calibration targets and controlled acquisition environments, which may not always be practical.

Alternative approaches have explored the use of natural or commonly occurring objects as calibration targets, in particular simple parametric surfaces such as spheres and cylinders. Among these, surfaces of revolution (SOR)~\cite{liu2019realtime} have attracted significant attention due to their symmetry properties and the feasibility of extracting their silhouettes in images. As early as \cite{Maybank1992}, the role of the absolute conic was established for self-calibration, offering methods that rely on invariants under projective transformations.

\subsection{Calibration from Spheres and Surfaces of Revolution}

The use of spheres as calibration objects has been extensively studied. \cite{Zhang2007} proposed a method to calibrate cameras directly from images of spheres by relating the dual images of spheres to the dual image of the absolute conic (DIAC). Their approach exploits the pole-polar relationship with respect to the image of the absolute conic, leading to a linear solution for estimating the intrinsic parameters from a minimum of three spheres. Earlier work by \cite{teramoto2002calibration} also utilized spheres to extract constraints via the absolute conic but required nonlinear optimization heavily dependent on initial estimates.

The robustness of sphere-based calibration lies in the ease of detecting spherical silhouettes under various viewpoints and partial occlusions, as emphasized by \cite{Zhang2007}. Spheres ensure visibility of their occluding contours from any vantage point, simplifying the extraction of reliable image features. However, the main limitation is the requirement for multiple distinct spheres or multiple views to provide enough independent constraints for a full calibration.

The broader family of SORs extends these ideas to other objects such as cylinders or cones. The inherent rotational symmetry provides strong geometric constraints that can be exploited for calibration, as demonstrated by \cite{Li2024}. These methods typically derive constraints from tangent cones or symmetry axes to compute the camera pose or intrinsic parameters.

\subsection{Pose Estimation from Cylindrical Features}

In parallel to calibration, significant work has addressed pose estimation using cylindrical primitives. Cylinders are abundant in industrial and urban environments, making them practical targets for vision-based localization. \cite{gummeson22pose} proposed a suite of minimal solvers to estimate camera pose from the silhouettes of multiple cylinders, even in the presence of unknown focal lengths. Their approach demonstrated robustness in both synthetic and real datasets, outperforming earlier methods, especially when combined with bundle adjustment.

The strength of cylinder-based methods lies in their ability to operate with few correspondences, sometimes even from a single frame, as the tangency condition of the image lines to the cylinder silhouette provides strong constraints. \cite{gummeson22pose} also demonstrated that their solvers could be integrated into RANSAC frameworks to robustly handle outliers, making them suitable for practical applications where object detections may be noisy or incomplete.

\subsection{Ellipsoid and Ellipse-Based Pose Estimation}

A growing body of literature has explored the use of ellipsoidal models for object-based pose estimation. Ellipsoids provide a more general parametric representation that encompasses spheres and spheroids as special cases. The Perspective-1-Ellipsoid (P1E) problem, as formalized by \cite{Gaudilliere2023}, investigates the estimation of camera pose from a single ellipse-ellipsoid correspondence.

Most traditional approaches employ the projective relationship between quadrics and conics, where a 3D ellipsoid projects to a 2D ellipse via the camera projection matrix. However, the general projective framework often fails to distinguish ellipsoids from other quadrics in a practical and robust manner. This ambiguity, combined with over-parameterization issues in the projection matrix, led \cite{Gaudilliere2023} to propose an ellipsoid-specific Euclidean formalism. This formulation decouples orientation and position estimation and reduces the pose estimation problem to a one-degree-of-freedom (1DoF) problem, enabling closed-form solutions under certain configurations.

Their approach demonstrates practical benefits, especially when few correspondences are available. Notably, it offers constructive solutions for both position-from-orientation and orientation-from-position scenarios, depending on available priors such as IMU-based orientation estimates or vanishing point detections. The sensitivity of the solution to ellipse detection noise has also been extensively analyzed, revealing that small perturbations in ellipse parameters can significantly affect orientation estimates \cite{gaudilliere2017perspective, gaudilliere2017perspective}. Ongoing research aims to enhance robustness by incorporating constraints on eigenvalue multiplicities and leveraging learning-based detection of ellipses to improve stability \cite{Zins2022}.

\subsection{Learning-Based Object-Aware Pose Estimation}

The advent of deep learning has introduced new possibilities for object-based pose estimation. \cite{Zins2022} introduced a 3D-aware ellipse prediction framework where neural networks are trained to directly regress ellipse parameters from 2D images, effectively bridging object detection with geometric pose computation. Their method relies on building a lightweight ellipsoid cloud model of the scene, where each object is approximated by a single ellipsoid described by only nine parameters.

In their pipeline, detected 2D ellipses are matched to their corresponding 3D ellipsoids, and the camera pose is then inferred by aligning projection and back-projection cones. This alignment translates the problem into one of cone intersection, echoing principles previously introduced by \cite{Gaudilliere2023}. The authors demonstrate that their method is robust to occlusions and viewpoint variations, provided accurate ellipse detections are obtained. They further propose hybrid pipelines combining deep learning-based detection with analytical geometric solvers, enabling real-time pose estimation with minimal scene-specific training data.

Such object-based methods offer clear advantages over classical keypoint-based approaches, especially in textureless or low-texture environments where sparse feature matching may fail. Moreover, these methods are particularly well-suited for robotic applications where semantic object understanding can be coupled with pose estimation.

\subsection{Hybrid and Multi-Object Approaches}

Beyond single-object solutions, several works have proposed combining multiple object observations or multiple geometric primitives to improve robustness and accuracy. For instance, \cite{Rubino2018} demonstrated the reconstruction of quadrics from multiple views using three calibrated cameras. While their method requires nonlinear post-processing to enforce ellipsoid constraints, it highlights the potential of multi-view integration.

In practice, multiple ellipse-ellipsoid correspondences can be exploited jointly to compute intersections of camera loci, as illustrated by \cite{gaudilliere2017perspective, gaudilliere2017perspective} using datasets such as T-LESS \cite{Sturm2004}. In such scenarios, minimal solvers provide hypotheses that are subsequently refined using RANSAC-based consensus schemes, achieving highly accurate localization even with few observations.

The integration of such hybrid strategies into SLAM and object-based structure-from-motion (SfM) frameworks has been investigated by \cite{Garg2020}, where semantic priors are incorporated into the optimization, enhancing robustness in cluttered or dynamic environments.

\subsection{Summary and Positioning of Our Work}

The body of research reviewed above highlights the diversity of approaches for camera calibration and pose estimation using higher-level geometric primitives. Spheres, cylinders, and ellipsoids provide rich sources of constraints that can compensate for the limitations of point-based methods, particularly in scenarios where artificial calibration targets or detailed textured models are impractical.

Our work builds upon these foundations by focusing specifically on infinite-length cylinders. While spheres and ellipsoids benefit from well-established dual conic and quadric relationships, and ellipsoids offer parameter-efficient models for general object shapes, infinite cylinders introduce distinctive constraints due to their translational invariance along the principal axis. This property allows us to formulate the pose estimation problem using a reduced set of parameters while avoiding certain degeneracies inherent in bounded primitives. Moreover, cylinders are pervasive in real-world scenes, from industrial environments to urban infrastructures, making them highly relevant for practical applications.

By drawing inspiration from both the analytical formulations presented for ellipsoids \cite{Gaudilliere2023} and the minimal solver frameworks developed for cylinders \cite{gummeson22pose}, we propose a novel method that exploits the algebraic structure of infinite cylinders within a consistent Euclidean framework. Our approach combines the strengths of prior work while addressing specific challenges unique to the cylinder case, such as axis alignment ambiguities and pose observability under limited viewpoint diversity.

In summary, our contribution extends the ongoing evolution of object-based camera calibration and pose estimation, offering a complementary perspective that leverages the advantageous properties of cylindrical geometries within a unified geometric formalism.

\begin{table}[tb]
\begin{center}
\begin{tabular}{@{}lllll@{}}
\toprule
                & Guo\cite{guo23vanishing}      & Ding \cite{CalibrationFreePlacedRightCylinder}        & Gummeson   \cite{gummeson22pose,gummeson24relative}           & Ours                    \\ \midrule
Calibration     & $f$ \mehmark & $\KK$ \cmark&  $f$  \mehmark             & $\KK $     \cmark       \\
Pose estimation &  $\RR$ \mehmark & \xmark     & $\RR,\mathbf{t}$ \cmark & $\RR,\mathbf{t}$ \cmark\\
\bottomrule
\end{tabular}
\end{center}
\caption{Comparison of camera estimation methods and their outputs for calibration and pose estimation. (\xmark: not addressed, \mehmark: partial output, \cmark: fully solved). Our method uniquely solves both tasks and provides all parameters.}
\label{tab:soa}
\end{table}
